\itemhead{第2章 小结与复习 · 第5题（教材p.60）}{XJ-B1-C02-SU-P060-Y1427}
\label{XJ-B1-C02-SU-P060-Y1427}
{\footnotesize 湘教版教材p.60 / PDF第67页 / 第2章 小结与复习 / 第5题}
\field{原题}{不等式可以刻画现实世界中的数量关系。尝试从实际问题中抽象出一元二次不等式模型，并进行求解，进而解决简单的实际问题。}
\field{参考答案}{示例：价格 $50+10x$、销量 $200-10x$，收入至少12600转为 $(50+10x)(200-10x)\ge12600$，解出后还需与 $x$ 的实际正整数范围取交集。}
{\small \field{关联知识}{I04 一元二次不等式与参数分类（教材帮p.64、65、66 / PDF73、74、75）；F03 函数求值、解析式与函数建模（教材帮p.80、87、88、89、90、131 / PDF89、96、97、98、99、140）}
\field{方法与易错点}{化标准式、查系数、求根或判别式、按图定号，最后与实际范围取交集。 易错：套“内小外大”却忘开口；丢重根或降次；漏整数产量限制。}
\field{答案依据}{答案册未提供本栏目的独立对应答案；依据原题、教材概念补写参考回应。开放性总结不设唯一数值答案。}}
